\documentclass[11pt]{article}
\usepackage{amsmath,amssymb,amsthm,mathtools}
\usepackage[margin=1in]{geometry}
\usepackage{booktabs}
\usepackage{enumitem}

\newtheorem{theorem}{Theorem}
\newtheorem{lemma}{Lemma}
\newtheorem{corollary}{Corollary}
\newtheorem{proposition}{Proposition}
\newtheorem{definition}{Definition}
\newtheorem{remark}{Remark}

\DeclareMathOperator{\Ran}{Ran}
\DeclareMathOperator{\Cost}{Cost}
\DeclareMathOperator{\Ker}{Ker}
\DeclareMathOperator{\tr}{tr}

\title{Step 53: Navier Adequacy Demonstration\\
\large Energy Ledgers, Blow-Up Probes, and the Adequacy Residual $\Xi$}
\author{RATCHET anti-localization workstream}
\date{}

\begin{document}
\maketitle

\section{Purpose and status}

This note is not a proof of Navier--Stokes regularity and it is not a numerical Six Birds simulation.  It is a symbolic adequacy demonstration: it translates the formed-layer membrane theory into the Navier-facing language of native energy/enstrophy probes versus layer-dissolving concentration probes.

The main point is that classical energy accounting can be correct and still be inadequate for pointwise or scale-local blow-up control.  In the membrane language, the missing quantity is the adequacy residual
\[
\Xi_C(D\mid L),
\]
which measures the part of a layer-dissolving probe family $D$ invisible to the declared native probe family $L$ under the audit geometry $C$.

\section{General adequacy setup}

Let $E$ be a Hilbert carrier, $C\succeq0$ an audit operator, and work on the legal quotient
\[
E_C=(\ker C)^\perp,
\qquad C_0=C|_{E_C}>0.
\]
Let
\[
L_0:E_C\to Y,
\qquad D_0:E_C\to Z
\]
be, respectively, the native probe family and a layer-dissolving probe family.  Define
\[
K_{LL}=L_0C_0^{-1}L_0^*,
\qquad
K_{DL}=D_0C_0^{-1}L_0^*,
\qquad
K_{DD}=D_0C_0^{-1}D_0^*.
\]
The adequacy residual is
\[
\boxed{
\Xi_C(D\mid L)=K_{DD}-K_{DL}K_{LL}^{\dagger}K_{LD}.
}
\]
It is the Schur residual of the dissolving currency after conditioning on native currency.

\begin{theorem}[Adequacy promotion]
If
\[
K_{LL}\preceq \Theta_Y,
\qquad
\Xi_C(D\mid L)\preceq \Xi_Z,
\]
then, with $A_*=K_{DL}K_{LL}^{\dagger}$,
\[
\boxed{
K_{DD}\preceq A_*\Theta_YA_*^*+\Xi_Z.
}
\]
In particular, a native membrane controls layer-dissolving probes only after adequacy is paid by $\Xi_Z$.
\end{theorem}

\begin{proof}
By definition,
\[
K_{DD}=K_{DL}K_{LL}^{\dagger}K_{LD}+\Xi_C(D\mid L)=A_*K_{LL}A_*^*+\Xi_C(D\mid L).
\]
Using $K_{LL}\preceq\Theta_Y$ and $\Xi_C(D\mid L)\preceq\Xi_Z$ gives the claim.
\end{proof}

\section{A Sobolev model for localized Navier probes}

Work in a local Euclidean chart or on a periodic box in dimension $d$.  Let the audit be the Sobolev energy
\[
\|u\|_{H^s}^2=\int (1+|\xi|^2)^s|\widehat u(\xi)|^2\,d\xi.
\]
Let $\varphi\in C_c^\infty(\mathbb R^d)$ be a normalized bump, and define a scale-$r$ localized derivative probe of order $m=|\alpha|$ by
\[
D_{x,r}^{\alpha}u
=
\int r^{-d}\varphi\!\left(\frac{y-x}{r}\right)\partial^\alpha u(y)\,dy.
\]
Up to harmless constants, its capacity under the $H^s$ audit is
\[
\operatorname{Cap}_s(D_{x,r}^{\alpha})
=
\int |\xi|^{2m}|\widehat\varphi(r\xi)|^2(1+|\xi|^2)^{-s}\,d\xi.
\]

\begin{theorem}[Localized Sobolev capacity scaling]
For a localized derivative probe of order $m$ in dimension $d$,
\[
\operatorname{Cap}_s(D_{x,r}^{\alpha})
\begin{cases}
\asymp r^{2(s-m)-d}, & s<d/2+m,\\
\asymp \log(1/r), & s=d/2+m,\\
O(1), & s>d/2+m,
\end{cases}
\]
as $r\downarrow0$, provided the bump is nondegenerate at low frequency.  Consequently, the point-scale derivative readout is bounded on $H^s$ only when
\[
\boxed{s>d/2+m.}
\]
\end{theorem}

\begin{proof}
For small $r$, high frequencies dominate when $s\le d/2+m$.  Substitute $\eta=r\xi$ in the principal high-frequency integral:
\[
\int |\xi|^{2m}(1+|\xi|^2)^{-s}|\widehat\varphi(r\xi)|^2\,d\xi
\sim
r^{2(s-m)-d}\int |\eta|^{2m-2s}|\widehat\varphi(\eta)|^2\,d\eta.
\]
The exponent $2(s-m)-d$ gives the power law.  The borderline exponent is zero and produces logarithmic divergence.  If $s>d/2+m$, the point evaluation derivative is a bounded functional on $H^s$ by the same Fourier integral criterion, giving a finite limit.
\end{proof}

\section{Navier-facing corollaries in dimension three}

For Navier--Stokes in three dimensions, a pointwise velocity probe corresponds to $m=0$, while a pointwise gradient, strain, or vorticity probe corresponds to $m=1$.

\begin{corollary}[Energy/enstrophy blind spots]
In dimension $d=3$:
\begin{enumerate}[label=(\roman*)]
\item The $L^2$ energy audit $s=0$ does not control pointwise velocity probes; scale-$r$ capacity diverges like $r^{-3}$.
\item The enstrophy/dissipation audit $s=1$ does not control pointwise velocity probes; scale-$r$ capacity diverges like $r^{-1}$.
\item The enstrophy/dissipation audit $s=1$ does not control pointwise gradient, vorticity, or strain probes; scale-$r$ capacity diverges like $r^{-3}$.
\item Even an $H^2$ audit does not control pointwise gradient/vorticity probes; the scale-$r$ gradient capacity diverges like $r^{-1}$.
\item Pointwise gradient/vorticity probes require $s>5/2$.
\end{enumerate}
\end{corollary}

\begin{proof}
Apply the previous theorem with $d=3$ and $m=0$ or $m=1$.  The threshold is $s>3/2$ for velocity and $s>5/2$ for gradients.
\end{proof}

\section{Interpreting $\Xi$ for Navier}

Let $L_j$ denote the declared native probe family at scale $j$: shell energies, filtered velocity observations, enstrophy-like measurements, or other probes admitted by the formed closure.  Let $D_j$ denote the layer-dissolving probes required by a regularity claim: pointwise or scale-local velocity, vorticity, strain, or blow-up-rate readouts.

The adequacy residual
\[
\boxed{
\Xi_{C_j}(D_j\mid L_j)
}
\]
measures the portion of the blow-up readout that is invisible to the native energy/enstrophy/shell ledger.

\begin{proposition}[Energy estimates are not adequacy]
A bounded native energy/enstrophy membrane
\[
K_{L_jL_j}\preceq\Theta_j
\]
does not imply a bounded blow-up-probe membrane unless
\[
\Xi_{C_j}(D_j\mid L_j)
\]
is also bounded, summable, or collapsed by a strict extension.  In particular, if the layer-dissolving probes include pointwise gradient/vorticity probes and the audit regularity satisfies $s\le5/2$ in $d=3$, then the instantaneous Sobolev model gives divergent capacity at vanishing scale.
\end{proposition}

\begin{proof}
The promotion theorem gives
\[
K_{D_jD_j}\preceq A_j\Theta_jA_j^*+\Xi_{C_j}(D_j\mid L_j).
\]
Thus the native membrane contributes only through the factorable part $A_j\Theta_jA_j^*$.  The remaining layer-dissolving currency is exactly $\Xi$.  The Sobolev scaling theorem shows that, for pointwise gradient/vorticity probes under insufficient regularity, the dissolving capacity diverges as scale $r\downarrow0$.  Hence bounded energy/enstrophy accounting alone cannot be an adequacy proof.
\end{proof}

\section{Strict-extension alternatives}

The symbolic Navier membrane route has several lawful alternatives:
\begin{enumerate}[label=(\arabic*)]
\item \textbf{Coercive audit extension.} Add a lawful audit source giving sufficient regularity or scale-local coercivity, e.g. a mechanism that effectively prices $H^{5/2+\varepsilon}$-strength gradient readouts.
\item \textbf{Native-family extension.} Add the relevant localized vorticity/strain probes to the native family, then prove their predictive family currency remains bounded.
\item \textbf{Parabolic smoothing record.} Use positive-time smoothing only with an audited bridge to the time at which blow-up is claimed; smoothing at fixed positive time does not by itself control instantaneous singularity formation.
\item \textbf{Nonclaim/gating.} Exclude pointwise blow-up probes from the claim.  This may yield a scoped membrane, but it is not a Navier regularity theorem.
\end{enumerate}

Each alternative must be recorded as a strict all-six extension.  Repeating the same energy/enstrophy ledger cannot reduce $\Xi$.

\section{Navier obligation ledger}

A Navier-facing formed-layer membrane proof must provide:
\begin{enumerate}[label=(O\arabic*)]
\item the exact carrier $E_j$ and legal quotient, including divergence-free and boundary/support records;
\item the positive audit $C_j$ and null-mode legality;
\item the native predictive probe family $L_j$;
\item the layer-dissolving probe family $D_j$ required by the regularity/continuation claim;
\item an adequacy residual bound $\Xi_{C_j}(D_j\mid L_j)\preceq\Xi_j$;
\item predictive transport and summable defects for both native and dissolving families;
\item all-six statuses for rewrite, feasibility, route, staging, packaging, and audit.
\end{enumerate}

\section{Conclusion}

The $\Xi$-language explains exactly why energy accounting is not enough for Navier-style regularity.  The issue is not that energy estimates are false; it is that the layer-dissolving probes required by blow-up exclusion need not factor through the native energy/enstrophy ledger.  The adequacy residual measures this blind spot.

Thus the Navier-facing membrane obligation is
\[
\boxed{
\text{native membrane}
+
\Xi\text{-adequacy for blow-up probes}
+
\text{predictive strict-extension control}
\Rightarrow
\text{no blow-up needles}.
}
\]
This note is a symbolic demonstration of the framework's generality, not an NS proof.

\end{document}
